\section{Protocol integration}
\label{sec:protocol-integration}

The code is wired end-to-end through Zip$^+$'s commit / prove / verify.
This section describes the verifier-side field and the challenge type
that make that possible. (An earlier draft of this document described
this as an open ``protocol-integration gap''; it is now closed, and the
content below supersedes it.)

\subsection{The obstruction}

The Zip$^+$ prover, at its combined-row / proximity step, computes
inner products of \code{Zt::CombR} slices against verifier-side
field elements. This requires a map
\[
  \mathrm{Zt::CombR} \;\longrightarrow\; F,
\]
i.e.\ $F : \mathrm{FromWithConfig}\langle \&\,\code{Zt::CombR}\rangle$.
For the IPRS / RAA codes \code{CombR} is a scalar \code{Int<N>} and
the map is the standard ``integer mod $p$'' lift into a prime field.
For this code, \code{CombR} is a degree-$16$ polynomial in $\Rlift$,
and no map into a \emph{prime} field $\F_p$ is natural: the encoder's
algebra is polynomial-shaped, so the verifier-side scalar must be too.

\subsection{The verifier-side field $F = \F_p[X]/\flift$}

The fix is to make $F$ itself polynomial-shaped. Define
\[
  F \;=\; \F_p[X]/\flift,
\]
the same quotient shape as $\Rlift$ but with prime-field coefficients
in place of $\Z$-coefficients. The map
\[
  \pi_p : \Rlift \longrightarrow F, \qquad
  \pi_p(\,\textstyle\sum a_i X^i\,) = \sum (a_i \bmod p)\, X^i
\]
is coefficient-wise mod-$p$ reduction. Because $f \equiv \flift
\pmod p$ for the integer-coefficient $\flift$, $\pi_p$ is a
\emph{ring homomorphism} $\Rlift \to F$ --- it commutes with addition,
multiplication, and the $\flift$ reduction --- so it preserves every
algebraic identity the proximity and eval-consistency checks rely on.

When $\flift \bmod p$ is irreducible over $\F_p$, $F$ is a genuine
field, the degree-$16$ extension $\F_{p^{16}}$. The implementation
is \code{Gf2\_16Ext<const P: u64>} in \code{ext\_field.rs}, a
\code{[u64; 16]} coefficient vector mod $p$. It carries
\code{Field::Modulus = Uint<16>} (width-matched to the inner
representation so the transcript's shared modulus / inner buffers
agree) and implements the \code{PrimeField} trait surface that
prove/verify actually exercises --- the trait name is a slight abuse,
but only the \code{Field} operations and config-handling methods are
used; field-specific methods (square roots, $(p-1)/2$) are stubbed.

\paragraph{Prime selection.}
$F$ is a field --- so that \code{Div} / \code{Inv} make sense --- only
when $\flift \bmod p$ is irreducible. \code{ext\_field.rs} ships a
Rabin-style irreducibility test (\code{f\_tilde\_is\_irreducible\_mod\_p})
and a prime scanner. The default prime is
\[
  \code{P\_16\_DEFAULT} \;=\; 2\,147\,482\,921,
\]
the largest prime below $2^{31}$ for which
$\flift = X^{16}+X^5+X^3+X^2+1$ is irreducible mod $p$. A non-ignored
test pins this choice; $P$ is a const-generic, so it can be swapped.

\subsection{The challenge type {\tt Bit4Poly16}}

The row-batching challenge $\code{Chal}$ and evaluation point
$\code{Pt}$ are also polynomial-shaped: \code{Bit4Poly16}, a
$16$-nibble polynomial with $4$-bit integer coefficients in
$\{0,\ldots,15\}$, drawn fresh via Fiat--Shamir. It is backed by a
single \code{i64} ($16 \times 4 = 64$ bits) and has an $8$-byte wire
encoding. The narrow $4$-bit coefficients keep the combined-row
magnitudes small: at $\code{num\_rows} = 1$, a combined-row
coefficient is a sum of $\le \code{batch}$ products of a binary
message bit and a $4$-bit challenge coefficient, hence bounded by
$\code{batch}\cdot 15$ --- which is why \code{BITS\_CR = 8} suffices
for the combined-row wire encoding (\autoref{sec:packed-cw}).

\subsection{The coherence identity}

With $F = \F_p[X]/\flift$ and $\pi_p$ a ring homomorphism, the
protocol's coherence identity holds by algebra: the prover's
combined row, the evaluation, and the tensor coordinates
$(q_0, q_1)$ all live in $F$, and
\[
  \langle \text{combined\_row}, q_1 \rangle_F
  \;=\;
  \langle \text{coeffs}, b \rangle_F
\]
because every quantity is the $\pi_p$-image of the corresponding
$\Rlift$-valued quantity the prover computed, and $\pi_p$ commutes
with the inner product. The transcript serializes the combined-row
vector $b$ as polynomial-valued (packed) entries rather than scalar
field elements.

\subsection{Soundness note}

The proximity argument for the lifted code in $F = \F_p[X]/\flift$ is
not given a formal proof here. The structural ingredients are: the
lifted code reduces mod $2$ to a genuine $\GF(2^{16})$ Reed--Solomon
code (\autoref{sec:lift}), and the Schwartz--Zippel-style argument
the Zip$^+$ proximity test uses has a clean analogue in the extension
$\F_{p^{16}}$ when $\flift \bmod p$ is irreducible. A complete
analysis is left as follow-up (\autoref{sec:future-work}).
